% Daftar Pustaka (gaya APA edisi ke-7)
\babtp{DAFTAR PUSTAKA}
\begin{pustaka}
AlQahtani, S. A., \& Alhassany, M. (2013). Comparing different LTE scheduling schemes. \textit{2013 9th International Wireless Communications and Mobile Computing Conference (IWCMC)}, 264--269.

Andrews, M. (2007). A survey of scheduling theory in wireless data networks. \textit{Wireless Communications, 143}, 1--17.

Capozzi, F., Piro, G., Grieco, L. A., Boggia, G., \& Camarda, P. (2013). Downlink packet scheduling in LTE cellular networks: Key design issues and a survey. \textit{IEEE Communications Surveys \& Tutorials, 15}(2), 678--700.

Dahlman, E., Parkvall, S., \& Sk\"{o}ld, J. (2011). \textit{4G: LTE/LTE-Advanced for mobile broadband}. Academic Press.

Ghosh, A., Zhang, J., Andrews, J. G., \& Muhamed, R. (2010). \textit{Fundamentals of LTE}. Prentice Hall.

Jain, R. K. (1991). \textit{The art of computer systems performance analysis: Techniques for experimental design, measurement, simulation, and modeling}. John Wiley \& Sons.

ns-3 Consortium. (t.t.). \textit{ns-3 network simulator}. Diakses pada 21 Oktober 2026, dari \url{https://www.nsnam.org}

Xu, Y., \& Fischione, C. (2012). Real-time scheduling in LTE for smart grids. \textit{Proceedings of the 5th International Symposium on Communications, Control and Signal Processing}, 2--4.
\end{pustaka}
